\documentclass[conference]{IEEEtran}
\IEEEoverridecommandlockouts
\usepackage{cite}
\usepackage{amsmath,amssymb,amsfonts}
\usepackage{algorithmic}
\usepackage{graphicx}
\usepackage{textcomp}
\usepackage{xcolor}
\usepackage{booktabs}
\graphicspath{{figures/}}
\def\BibTeX{{\rm B\kern-.05em{\sc i\kern-.025em b}\kern-.08em
    T\kern-.1667em\lower.7ex\hbox{E}\kern-.125emX}}

\begin{document}

\title{Progress Report: Performance Comparison of Semantic Communication vs.\ Raw-Bit Transmission}

\author{
\IEEEauthorblockN{Jihyeok Lee}
\IEEEauthorblockA{
    \textit{Dept. of Computer Engineering} \\
    \textit{Chosun University} \\
    Gwangju, Republic of Korea \\
    dlwlgur02@chosun.ac.kr
}
}

\maketitle

% ──────────────────────────────────────────────
\begin{abstract}
This progress report summarizes the work completed to date on a performance
comparison between semantic communication and raw-bit (JPEG-compressed)
transmission for task-oriented mobile computing.
We established a complete simulation pipeline using CIFAR-10 image
classification as the downstream task.
A ResNet-18 classifier, serving as the baseline, achieved a test accuracy of
93.51\%.
Semantic encoders with latent dimensions of 64, 128, and 256 were trained
end-to-end and evaluated under AWGN and packet-loss channel conditions.
Initial results show that the semantic approach (dim\,=\,64, 64\,bytes) achieves
91.94\% accuracy at SNR\,=\,10\,dB---using only 2.1\% of the bandwidth required
by the lossless baseline---while outperforming JPEG compression by a factor of
approximately 15$\times$ in accuracy per bit.
\end{abstract}

% ──────────────────────────────────────────────
\section{Introduction}

The proliferation of AI-driven mobile applications demands efficient data
transmission for inference tasks over bandwidth-constrained wireless channels.
Conventional raw-bit transmission, even with source coding such as JPEG,
transmits pixel-level information regardless of the downstream task, leading to
poor utilization of limited channel resources.
Semantic communication addresses this issue by extracting and transmitting only
task-relevant features, thereby reducing the payload while preserving task
performance \cite{b1,b2}.

This project quantitatively compares semantic communication and JPEG-based
raw-bit transmission under realistic wireless channel conditions, focusing on
three metrics: Top-1 classification accuracy, accuracy per bit~(ApB), and
robustness to channel degradation (SNR variation and packet loss).

% ──────────────────────────────────────────────
\section{Problem Statement}

Given a bandwidth-constrained wireless channel, the transmitter must send a
representation of an image $x$ such that the receiver can correctly classify it.
The key tension is between the transmission size (bytes) and task accuracy.
For a transmission budget $B$ (bytes), we compare two approaches:

\begin{itemize}
    \item \textbf{Baseline:} JPEG compression at quality $q$, where image $x$
          is compressed to approximately $B$ bytes, transmitted, decompressed,
          and fed to a classifier.
    \item \textbf{Semantic:} A learned encoder $f_{\phi}(x) \in \mathbb{R}^{d}$
          extracts a $d$-dimensional latent vector, quantized to $d$ bytes via
          8-bit scalar quantization, transmitted through a noisy channel, and
          classified by a lightweight head $g_{\psi}$.
\end{itemize}

The channel is modeled as AWGN with varying SNR, Rayleigh fading, or
block-level packet loss.

% ──────────────────────────────────────────────
\section{Work Completed}

\subsection{Simulation Environment}

The full experimental pipeline was implemented in Python/PyTorch and
version-controlled on GitHub. The system runs on an Apple MPS (M-series GPU).
Key components are as follows.

\begin{itemize}
    \item \textbf{Dataset:} CIFAR-10 (50{,}000 train / 10{,}000 test,
          $32{\times}32$ RGB)
    \item \textbf{Channel simulators:} AWGN, Rayleigh fading, block-level
          packet loss
    \item \textbf{Quantizer:} 8-bit uniform scalar quantization
    \item \textbf{Baseline:} JPEG compression via PIL at quality levels
          $q \in \{1, 3, 5, 10, 20, 40, 60, 80, 95\}$
\end{itemize}

\subsection{Baseline: ResNet-18 Classifier}

A ResNet-18 classifier was adapted for CIFAR-10 ($3{\times}3$ first
convolution, stride~1, no max-pooling) and trained for 30 epochs using SGD
with cosine annealing (lr\,=\,0.1, weight decay\,=\,$5{\times}10^{-4}$,
momentum\,=\,0.9).
The model achieved a test accuracy of \textbf{93.51\%}, which serves as the
upper-bound reference for lossless transmission.

\subsection{Semantic Encoder Training}

Three SemanticEncoder models (latent dim $\in \{64, 128, 256\}$) were trained
end-to-end using the same ResNet-18 backbone with a linear projection head,
paired with a two-layer MLP classifier head.
Training used Adam (lr\,=\,$1{\times}10^{-3}$, cosine annealing, 30 epochs).
Results are summarized in Table~\ref{tab:semantic_train}.

\begin{table}[htbp]
\caption{Semantic Encoder Training Results}
\label{tab:semantic_train}
\begin{center}
\begin{tabular}{ccc}
\toprule
\textbf{Latent Dim} & \textbf{Transmission Size} & \textbf{Best Test Acc.} \\
\midrule
64  & 64\,B  & 92.52\% \\
128 & 128\,B & 92.25\% \\
256 & 256\,B & 92.42\% \\
\bottomrule
\end{tabular}
\end{center}
\end{table}

\subsection{Baseline Evaluation (JPEG)}

JPEG compression was evaluated at nine quality levels.
The ResNet-18 classifier (Section~III-B) was applied to decompressed images
without retraining. Results are shown in Table~\ref{tab:jpeg}.

\begin{table}[htbp]
\caption{JPEG Baseline Evaluation Results}
\label{tab:jpeg}
\begin{center}
\begin{tabular}{lccc}
\toprule
\textbf{Method} & \textbf{Avg.\ Bytes} & \textbf{Top-1 Acc.} & \textbf{ApB ($\times10^{-4}$)} \\
\midrule
Raw (lossless) & 3{,}072\,B & 93.51\% & 0.38 \\
JPEG $q$=10    & 699.5\,B   & 42.47\% & 0.76 \\
JPEG $q$=40    & 801.3\,B   & 72.63\% & 1.13 \\
JPEG $q$=60    & 853.8\,B   & 78.83\% & 1.15 \\
JPEG $q$=80    & 960.6\,B   & 85.36\% & 1.11 \\
JPEG $q$=95    & 1{,}300.8\,B & 91.70\% & 0.88 \\
\bottomrule
\end{tabular}
\end{center}
\end{table}

JPEG reaches its peak accuracy per bit at $q$=60 ($1.15{\times}10^{-4}$) but
requires over 850\,bytes for approximately 79\% accuracy.

\begin{figure}[htbp]
\centerline{\includegraphics[width=\columnwidth]{acc_vs_byte_budget.png}}
\caption{Top-1 accuracy vs.\ transmission size (bytes). Semantic encoders
(stars) achieve comparable accuracy to JPEG (dashed line) at a fraction of
the byte budget.}
\label{fig:acc_byte}
\end{figure}

\begin{figure}[htbp]
\centerline{\includegraphics[width=\columnwidth]{accuracy_per_bit.png}}
\caption{Accuracy per bit vs.\ transmission size. Semantic encoders
consistently outperform JPEG across all byte budgets.}
\label{fig:apb}
\end{figure}

\subsection{Semantic Evaluation}

Semantic encoders were evaluated under AWGN
($\text{SNR} \in \{0, 5, 10, 20\}$\,dB) and packet loss
(rate $\in \{1\%, 5\%, 10\%\}$).
Results are presented in Tables~\ref{tab:awgn}--\ref{tab:efficiency}.

\begin{table}[htbp]
\caption{Semantic Accuracy under AWGN Channel (CIFAR-10 Test Set)}
\label{tab:awgn}
\begin{center}
\begin{tabular}{ccccc}
\toprule
\textbf{Latent Dim} & \textbf{0\,dB} & \textbf{5\,dB} & \textbf{10\,dB} & \textbf{20\,dB} \\
\midrule
64  & 88.67\% & 91.32\% & 91.94\% & 92.52\% \\
128 & 90.49\% & 91.92\% & 92.14\% & 92.40\% \\
256 & 91.43\% & 92.16\% & 92.16\% & 92.37\% \\
\bottomrule
\end{tabular}
\end{center}
\end{table}

\begin{table}[htbp]
\caption{Semantic Accuracy under Packet Loss}
\label{tab:pktloss}
\begin{center}
\begin{tabular}{cccc}
\toprule
\textbf{Latent Dim} & \textbf{Loss 1\%} & \textbf{Loss 5\%} & \textbf{Loss 10\%} \\
\midrule
64  & 92.46\% & 92.24\% & 92.32\% \\
128 & 92.21\% & 92.32\% & 92.18\% \\
256 & 92.40\% & 92.41\% & 92.45\% \\
\bottomrule
\end{tabular}
\end{center}
\end{table}

\begin{table}[htbp]
\caption{Communication Efficiency at SNR = 10\,dB}
\label{tab:efficiency}
\begin{center}
\begin{tabular}{lccc}
\toprule
\textbf{Method} & \textbf{Bytes} & \textbf{Top-1 Acc.} & \textbf{ApB ($\times10^{-4}$)} \\
\midrule
JPEG $q$=60 (best ApB) & 853.8  & 78.83\% & 1.15 \\
JPEG $q$=95            & 1{,}300.8 & 91.70\% & 0.88 \\
\textbf{Semantic dim=64}  & \textbf{64}  & \textbf{91.94\%} & \textbf{17.96} \\
Semantic dim=128       & 128    & 92.14\% & 8.99 \\
Semantic dim=256       & 256    & 92.16\% & 4.50 \\
\bottomrule
\end{tabular}
\end{center}
\end{table}

\begin{figure}[htbp]
\centerline{\includegraphics[width=\columnwidth]{snr_robustness.png}}
\caption{Top-1 accuracy vs.\ SNR (dB) under AWGN channel. Higher latent
dimensions exhibit greater robustness at low SNR.}
\label{fig:snr}
\end{figure}

\begin{figure}[htbp]
\centerline{\includegraphics[width=\columnwidth]{packet_loss_robustness.png}}
\caption{Top-1 accuracy vs.\ packet loss rate. All semantic models maintain
accuracy above 92\% even at 10\% packet loss.}
\label{fig:pktloss}
\end{figure}

% ──────────────────────────────────────────────
\section{Preliminary Findings}

\textbf{(1) Dramatic bandwidth reduction.}
As shown in Fig.~\ref{fig:acc_byte}, Semantic dim=64 achieves 91.94\% accuracy
using only 64\,bytes---2.1\% of the 3{,}072-byte lossless baseline---while JPEG
requires 1{,}301\,bytes to reach comparable accuracy (91.70\%).

\textbf{(2) Communication efficiency.}
Fig.~\ref{fig:apb} shows that Semantic dim=64 achieves an accuracy per bit of
$17.96{\times}10^{-4}$, approximately \textbf{15.6$\times$ higher} than the
best JPEG configuration ($q$=60, $1.15{\times}10^{-4}$).

\textbf{(3) AWGN robustness.}
As illustrated in Fig.~\ref{fig:snr}, even at SNR\,=\,0\,dB, Semantic dim=256
maintains 91.43\% accuracy. The dim=64 model shows greater sensitivity to low
SNR (88.67\% at 0\,dB) owing to its compressed representation.

\textbf{(4) Packet loss robustness.}
Fig.~\ref{fig:pktloss} shows that all semantic models are remarkably robust to
packet loss. The dim=256 accuracy fluctuates by less than 0.05\% across
1\%--10\% loss rates, suggesting that the latent representation distributes
task-relevant information redundantly across dimensions.

% ──────────────────────────────────────────────
\section{Updated Project Schedule}

Table~\ref{tab:gantt} summarizes the updated Gantt chart.
No major deviations from the original plan occurred.
The full evaluation of Rayleigh fading was deferred to the next phase
(April 9--20) to prioritize the submission of this report.

\begin{table}[htbp]
\caption{Updated Project Schedule (Gantt Chart)}
\label{tab:gantt}
\begin{center}
\begin{tabular}{p{2.6cm} p{2.8cm} p{1.8cm}}
\toprule
\textbf{Period} & \textbf{Activity} & \textbf{Milestone} \\
\midrule
03/13 -- 03/17 & Environment setup, baseline training & --- \\
03/18 & Proposal submission & Proposal \checkmark \\
03/19 -- 04/07 & Semantic encoder dev., initial experiments & --- \\
04/08 & Progress Report submission & This report \checkmark \\
04/09 -- 04/20 & Extended experiments (Rayleigh, Tiny-ImageNet) & --- \\
04/21 -- 05/10 & Statistical repetition (5 runs) & --- \\
05/11 -- 05/20 & Final report writing, slides & --- \\
05/27 & Presentation & Slides + Demo \\
06/03 & Final report submission & IEEE format \\
\bottomrule
\end{tabular}
\end{center}
\end{table}

% ──────────────────────────────────────────────
\section{Remaining Work}

\begin{itemize}
    \item \textbf{Extended channel evaluation:} Full Rayleigh fading sweep
          across all latent dimensions.
    \item \textbf{Statistical rigor:} Repeat all experiments 5 times (seed
          sweep) and report mean $\pm$ std.
    \item \textbf{Tiny-ImageNet:} Extend dataset if time permits (originally
          planned as optional).
    \item \textbf{Vector quantization:} Compare K-means VQ against scalar
          quantization.
    \item \textbf{Final report:} IEEE conference format, 10$+$ references.
\end{itemize}

% ──────────────────────────────────────────────
\begin{thebibliography}{00}

\bibitem{b1}
Z. Qin, X. Tao, J. Lu, and G. Y. Li,
``Semantic communications: Principles and challenges,''
\textit{arXiv preprint arXiv:2212.00032}, 2021.

\bibitem{b2}
S. Sukhbaatar, A. Szlam, and J. Weston,
``Learning multiagent communication with backpropagation,''
\textit{arXiv preprint arXiv:1605.07736}, 2016.

\bibitem{b3}
E. Bourtsoulatze, D. B. Kurka, and D. G\"{u}nd\"{u}z,
``Deep joint source-channel coding for wireless image transmission,''
\textit{IEEE Trans. Cogn. Commun. Netw.}, vol.~5, no.~3, pp.~567--579, 2019.

\bibitem{b4}
K. He, X. Zhang, S. Ren, and J. Sun,
``Deep residual learning for image recognition,''
in \textit{Proc. IEEE CVPR}, 2016, pp.~770--778.

\end{thebibliography}

\end{document}
